\hypertarget{chapter-6-polymorphism-and-virtualization}{%
\section{CHAPTER 6: POLYMORPHISM AND
VIRTUALIZATION}\label{chapter-6-polymorphism-and-virtualization}}

\hypertarget{polymorphism-virtualization}{%
\section{POLYMORPHISM \&
VIRTUALIZATION}\label{polymorphism-virtualization}}

\hypertarget{deep-object-model-virtualization}{%
\section{DEEP OBJECT MODEL \&
VIRTUALIZATION}\label{deep-object-model-virtualization}}

Understanding the ``C++ Object Model'' distinguishes a user from a
master. This section explains what the compiler generates for your
classes.

\hypertarget{the-cost-of-polymorphism-vptr-vtable}{%
\subsubsection{2.5.1 The Cost of Polymorphism (vptr \&
vtable)}\label{the-cost-of-polymorphism-vptr-vtable}}

Every class with \emph{at least one} virtual function has a hidden
overhead.

\begin{enumerate}
\def\labelenumi{\arabic{enumi}.}
\tightlist
\item
  \textbf{vptr (Virtual Pointer)}: A hidden member added to the
  \emph{layout} of the object.
\item
  \textbf{vtable (Virtual Table)}: A static table of function pointers
  for that class.
\end{enumerate}

\begin{lstlisting}[language={C++}]
class Base {
    int data;
    virtual void func() {}
};
// sizeof(Base) = sizeof(int) + sizeof(void*) + padding
// On 64-bit: 4 bytes (int) + 4 bytes (padding) + 8 bytes (ptr) = 16 bytes
\end{lstlisting}

\begin{center}\rule{0.5\linewidth}{0.5pt}\end{center}

\hypertarget{professional-notes-polymorphism-virtual-mechanics}{%
\subsubsection{Professional Notes: Polymorphism \& Virtual
Mechanics}\label{professional-notes-polymorphism-virtual-mechanics}}

\hypertarget{polymorphism-and-virtual-destructors}{%
\paragraph{1. Polymorphism and Virtual
Destructors}\label{polymorphism-and-virtual-destructors}}

Always declare the destructor as \passthrough{\lstinline!virtual!} in a
polymorphic base class. * \textbf{The Danger}: If you delete a derived
object via a base class pointer without a virtual destructor, only the
base part is destroyed, leading to \textbf{Memory Leaks}.

\begin{lstlisting}[language={C++}]
class Base {
public:
    virtual ~Base() {} // Essential!
};
\end{lstlisting}

\hypertarget{pure-virtual-functions-and-abstract-classes}{%
\paragraph{2. Pure Virtual Functions and Abstract
Classes}\label{pure-virtual-functions-and-abstract-classes}}

A function with \passthrough{\lstinline!= 0!} is pure virtual. It forces
derived classes to provide an implementation. * \textbf{Abstract Class}:
A class with at least one pure virtual function cannot be instantiated.
* \textbf{Interface}: In C++, interfaces are typically classes with only
public pure virtual functions and a virtual destructor.

\hypertarget{virtual-functions-in-constructorsdestructors}{%
\paragraph{3. Virtual Functions in
Constructors/Destructors}\label{virtual-functions-in-constructorsdestructors}}

\textbf{Godhood Warning}: Never call virtual functions in constructors
or destructors. * \textbf{Reason}: During the base class constructor,
the derived class members haven't been initialized yet. The vtable still
points to the base class implementation. This is a common source of
bugs.

\hypertarget{the-override-and-final-keywords-c11}{%
\paragraph{\texorpdfstring{4. The \texttt{override} and \texttt{final}
Keywords
(C++11)}{4. The override and final Keywords (C++11)}}\label{the-override-and-final-keywords-c11}}

\begin{itemize}
\tightlist
\item
  \textbf{\passthrough{\lstinline!override!}}: Ensures the function
  actually overrides a base class virtual function. Catches signature
  mismatches at compile time.
\item
  \textbf{\passthrough{\lstinline!final!}}: Prevents further overriding
  or inheritance.
\end{itemize}

\begin{center}\rule{0.5\linewidth}{0.5pt}\end{center}

\hypertarget{multiple-inheritance-thunks}{%
\subsubsection{2.5.2 Multiple Inheritance \&
Thunks}\label{multiple-inheritance-thunks}}

When inheriting from multiple classes, pointer arithmetic gets tricky.

\begin{lstlisting}[language={C++}]
class A { int a; virtual void f() {} };
class B { int b; virtual void g() {} };
class C : public A, public B { int c; };

C obj;
A* pa = &obj; // Points to start of obj
B* pb = &obj; // Points to obj + sizeof(A) !!
\end{lstlisting}

\begin{itemize}
\tightlist
\item
  \textbf{Thunk}: A small piece of assembly code generated by the
  compiler to adjust the \passthrough{\lstinline!this!} pointer when
  calling a virtual function from a base class pointer that isn't at
  offset 0.
\end{itemize}

\hypertarget{virtual-inheritance-the-diamond-problem}{%
\subsubsection{2.5.3 Virtual Inheritance (The Diamond
Problem)}\label{virtual-inheritance-the-diamond-problem}}

\begin{lstlisting}[language={C++}]
class Top { int t; };
class Left : virtual public Top { int l; };
class Right : virtual public Top { int r; };
class Bottom : public Left, public Right { int b; };
\end{lstlisting}

To solve the Diamond Problem (Top appearing twice),
\passthrough{\lstinline!virtual!} inheritance ensures
\passthrough{\lstinline!Top!} is shared. * \textbf{Cost}:
\passthrough{\lstinline!Left!} and \passthrough{\lstinline!Right!} now
contain a \textbf{vbptr} (Virtual Base Pointer) pointing to the shared
\passthrough{\lstinline!Top!} instance. Accessing members of
\passthrough{\lstinline!Top!} becomes slower (indirection).

\hypertarget{alignment-padding-rules}{%
\subsubsection{2.5.4 Alignment \& Padding
Rules}\label{alignment-padding-rules}}

CPU reads are efficient at specific addresses (multiples of 4 or 8).
Compilers insert ``padding bytes''.

\textbf{Rule}: A member of size \(N\) must sit at an offset divisible by
\(N\).

\begin{lstlisting}[language={C++}]
struct Mixed {
    char a;     // 1 byte
                // 3 bytes PADDING
    int b;      // 4 bytes
    short c;    // 2 bytes
                // 6 bytes PADDING (to align structure size to 8)
};
// sizeof(Mixed) = 16 (on 64-bit)
\end{lstlisting}

\textbf{Optimization}: Sort members by size (Largest to Smallest) to
minimize padding.

\begin{center}\rule{0.5\linewidth}{0.5pt}\end{center}

\hypertarget{professional-notes-chapter-25-polymorphism}{%
\section{Professional Notes: Chapter 25:
Polymorphism}\label{professional-notes-chapter-25-polymorphism}}

Section 25.1: Dene polymorphic classes Section 25.2: Safe downcasting
Section 25.3: Polymorphism \& Destructors

\hypertarget{professional-notes-chapter-38-virtual-member-functions}{%
\section{Professional Notes: Chapter 38: Virtual Member
Functions}\label{professional-notes-chapter-38-virtual-member-functions}}

Section 38.1: Final virtual functions Section 38.2: Using override with
virtual in C++11 and later Section 38.3: Virtual vs non-virtual member
functions Section 38.4: Behaviour of virtual functions in constructors
and destructors Section 38.5: Pure virtual functions

\hypertarget{professional-notes-chapter-40-special-member-functions}{%
\section{Professional Notes: Chapter 40: Special Member
Functions}\label{professional-notes-chapter-40-special-member-functions}}

Section 40.1: Default Constructor Section 40.2: Destructor Section 40.3:
Copy and swap Section 40.4: Implicit Move and Copy

\hypertarget{professional-notes-chapter-25-polymorphism-1}{%
\section{Professional Notes: Chapter 25:
Polymorphism}\label{professional-notes-chapter-25-polymorphism-1}}

Section 25.1: Dene polymorphic classes The typical example is an
abstract shape class, that can then be derived into squares, circles,
and other concrete shapes. The parent class: Let's start with the
polymorphic class: class Shape \{ public: virtual
\textasciitilde Shape() = default; virtual double get\_surface() const =
0; virtual void describe\_object() const \{ std::cout
\textless\textless{} ``this is a shape'' \textless\textless{} std::endl;
\} double get\_doubled\_surface() const \{ return 2 * get\_surface(); \}
\}; How to read this denition ? You can dene polymorphic behavior by
introduced member functions with the keyword virtual. Here
get\_surface() and describe\_object() will obviously be implemented
dierently for a square than for a circle. When the function is invoked
on an object, function corresponding to the real class of the object
will be determined at runtime. It makes no sense to dene get\_surface()
for an abstract shape. This is why the function is followed by = 0. This
means that the function is pure virtual function. A polymorphic class
should always dene a virtual destructor. You may dene non virtual member
functions. When these function will be invoked for an object, the
function will be chosen depending on the class used at compile-time.
Here get\_double\_surface() is dened in this way. A class that contains
at least one pure virtual function is an abstract class. Abstract
classes cannot be instantiated. You may only have pointers or references
of an abstract class type. Derived classes Once a polymorphic base class
is dened you can derive it. For example: class Square : public Shape \{
Point top\_left; double side\_length; public: Square (const Point\&
top\_left, double side) : top\_left(top\_left),
side\_length(side\_length) \{\} double get\_surface() override \{ return
side\_length * side\_length; \} void describe\_object() override \{
std::cout \textless\textless{} ``this is a square starting at''
\textless\textless{} top\_left.x \textless\textless{} ``,''
\textless\textless{} top\_left.y \textless\textless{} " with a length of
" \textless\textless{} side\_length \textless\textless{} std::endl; \}
\}; Some explanations: You can dene or override any of the virtual
functions of the parent class. The fact that a function was virtual in
the parent class makes it virtual in the derived class. No need to tell
the compiler the keyword virtual again. But it's recommended to add the
keyword override at the end of the function declaration, in order to
prevent subtle bugs caused by unnoticed variations in the function
signature. If all the pure virtual functions of the parent class are
dened you can instantiate objects for this class, else it will also
become an abstract class. You are not obliged to override all the
virtual functions. You can keep the version of the parent if it suits
your need. Example of instantiation int main() \{ Square
square(Point(10.0, 0.0), 6); // we know it's a square, the compiler also
square.describe\_object(); std::cout \textless\textless{} ``Surface:''
\textless\textless{} square.get\_surface() \textless\textless{}
std::endl; Circle circle(Point(0.0, 0.0), 5); Shape \emph{ps = nullptr;
// we don't know yet the real type of the object ps = \&circle; // it's
a circle, but it could as well be a square
ps-\textgreater describe\_object(); std::cout \textless\textless{}
``Surface:'' \textless\textless{} ps-\textgreater get\_surface()
\textless\textless{} std::endl; \} Section 25.2: Safe downcasting
Suppose that you have a pointer to an object of a polymorphic class:
Shape }ps; // see example on defining a polymorphic class ps =
get\_a\_new\_random\_shape(); // if you don't have such a function yet,
you // could just write ps = new Square(0.0,0.0, 5); a downcast would be
to cast from a general polymorphic Shape down to one of its derived and
more specic shape like Square or Circle. Why to downcast ? Most of the
time, you would not need to know which is the real type of the object,
as the virtual functions allow you to manipulate your object
independently of its type: std::cout \textless\textless{} ``Surface:''
\textless\textless{} ps-\textgreater get\_surface() \textless\textless{}
std::endl; If you don't need any downcast, your design would be perfect.
However, you may need sometimes to downcast. A typical example is when
you want to invoke a non virtual function that exist only for the child
class. Consider for example circles. Only circles have a diameter. So
the class would be dened as : class Circle: public Shape \{ // for
Shape, see example on defining a polymorphic class Point center; double
radius; public:  Circle (const Point\& center, double radius) :
center(center), radius(radius) \{\} double get\_surface() const override
\{ return r * r * M\_PI; \} // this is only for circles. Makes no sense
for other shapes double get\_diameter() const \{ return 2 * r; \} \};
The get\_diameter() member function only exist for circles. It was not
dened for a Shape object: Shape* ps = get\_any\_shape();
ps-\textgreater get\_diameter(); // OUCH !!! Compilation error How to
downcast ? If you'd know for sure that ps points to a circle you could
opt for a static\_cast: std::cout \textless\textless{} ``Diameter:''
\textless\textless{}
static\_cast\textless Circle\emph{\textgreater(ps)-\textgreater get\_diameter()
\textless\textless{} std::endl; This will do the trick. But it's very
risky: if ps appears to by anything else than a Circle the behavior of
your code will be undened. So rather than playing Russian roulette, you
should safely use a dynamic\_cast. This is specically for polymorphic
classes : int main() \{ Circle circle(Point(0.0, 0.0), 10); Shape
\&shape = circle; std::cout \textless\textless{} ``The shape has a
surface of'' \textless\textless{} shape.get\_surface()
\textless\textless{} std::endl; //shape.get\_diameter(); // OUCH !!!
Compilation error Circle }pc =
dynamic\_cast\textless Circle\emph{\textgreater(\&shape); // will be
nullptr if ps wasn't a circle if (pc) std::cout \textless\textless{}
``The shape is a circle of diameter'' \textless\textless{}
pc-\textgreater get\_diameter() \textless\textless{} std::endl; else
std::cout \textless\textless{} ``The shape isn't a circle !''
\textless\textless{} std::endl; \} Note that dynamic\_cast is not
possible on a class that is not polymorphic. You'd need at least one
virtual function in the class or its parents to be able to use it.
Section 25.3: Polymorphism \& Destructors If a class is intended to be
used polymorphically, with derived instances being stored as base
pointers/references, its base class' destructor should be either virtual
or protected. In the former case, this will cause object destruction to
check the vtable, automatically calling the correct destructor based on
the dynamic type. In the latter case, destroying the object through a
base class pointer/reference is disabled, and the object can only be
deleted when explicitly treated as its actual type. struct
VirtualDestructor \{ virtual \textasciitilde VirtualDestructor() =
default; \}; struct VirtualDerived : VirtualDestructor \{\}; struct
ProtectedDestructor \{ protected: \textasciitilde ProtectedDestructor()
= default; \}; struct ProtectedDerived : ProtectedDestructor \{
\textasciitilde ProtectedDerived() = default; \}; // \ldots{}
VirtualDestructor} vd = new VirtualDerived; delete vd; // Looks up
VirtualDestructor::\textasciitilde VirtualDestructor() in vtable, sees
it's // VirtualDerived::\textasciitilde VirtualDerived(), calls that.
ProtectedDestructor* pd = new ProtectedDerived; delete pd; // Error:
ProtectedDestructor::\textasciitilde ProtectedDestructor() is protected.
delete static\_cast\textless ProtectedDerived*\textgreater(pd); // Good.
Both of these practices guarantee that the derived class' destructor
will always be called on derived class instances, preventing memory
leaks.

\hypertarget{professional-notes-chapter-38-virtual-member-functions-1}{%
\section{Professional Notes: Chapter 38: Virtual Member
Functions}\label{professional-notes-chapter-38-virtual-member-functions-1}}

Section 38.1: Final virtual functions C++11 introduced final specier
which forbids method overriding if appeared in method signature: class
Base \{ public: virtual void foo() \{ std::cout \textless\textless{}
``Base::Foo\n''; \} \}; class Derived1 : public Base \{ public: //
Overriding Base::foo void foo() final \{ std::cout \textless\textless{}
``Derived1::Foo\n''; \} \}; class Derived2 : public Derived1 \{ public:
// Compilation error: cannot override final method virtual void foo() \{
std::cout \textless\textless{} ``Derived2::Foo\n''; \} \}; The specier
final can only be used with `virtual' member function and can't be
applied to non-virtual member functions Like final, there is also an
specier caller `override' which prevent overriding of virtual functions
in the derived class. The speciers override and final may be combined
together to have desired eect: class Derived1 : public Base \{ public:
void foo() final override \{ std::cout \textless\textless{}
``Derived1::Foo\n''; \} \}; Section 38.2: Using override with virtual in
C++11 and later The specier override has a special meaning in C++11
onwards, if appended at the end of function signature. This signies that
a function is Overriding the function present in base class \& The Base
class function is virtual There is no run time signicance of this
specier as is mainly meant as an indication for compilers The example
below will demonstrate the change in behaviour with our without using
override. Without override: \#include struct X \{ virtual void f() \{
std::cout \textless\textless{} ``X::f()\n''; \} \}; struct Y : X \{ //
Y::f() will not override X::f() because it has a different signature, //
but the compiler will accept the code (and silently ignore Y::f()).
virtual void f(int a) \{ std::cout \textless\textless{} a
\textless\textless{} ``\n''; \} \}; With override: \#include struct X \{
virtual void f() \{ std::cout \textless\textless{} ``X::f()\n''; \} \};
struct Y : X \{ // The compiler will alert you to the fact that Y::f()
does not // actually override anything. virtual void f(int a) override
\{ std::cout \textless\textless{} a \textless\textless{} ``\n''; \} \};
Note that override is not a keyword, but a special identier which only
may appear in function signatures. In all other contexts override still
may be used as an identier: void foo() \{ int override = 1; // OK. int
virtual = 2; // Compilation error: keywords can't be used as
identifiers. \} Section 38.3: Virtual vs non-virtual member functions
With virtual member functions: \#include struct X \{ virtual void f() \{
std::cout \textless\textless{} ``X::f()\n''; \} \}; struct Y : X \{ //
Specifying virtual again here is optional // because it can be inferred
from X::f(). virtual void f() \{ std::cout \textless\textless{}
``Y::f()\n''; \} \}; void call(X\& a) \{ a.f(); \} int main() \{ X x; Y
y; call(x); // outputs ``X::f()'' call(y); // outputs ``Y::f()'' \}
Without virtual member functions: \#include struct X \{ void f() \{
std::cout \textless\textless{} ``X::f()\n''; \} \}; struct Y : X \{ void
f() \{ std::cout \textless\textless{} ``Y::f()\n''; \} \}; void call(X\&
a) \{ a.f(); \} int main() \{ X x; Y y; call(x); // outputs ``X::f()''
call(y); // outputs ``X::f()'' \} Section 38.4: Behaviour of virtual
functions in constructors and destructors The behaviour of virtual
functions in constructors and destructors is often confusing when rst
encountered. \#include using namespace std; class base \{ public: base()
\{ f(``base constructor''); \} \textasciitilde base() \{ f(``base
destructor''); \} virtual const char* v() \{ return ``base::v()''; \}
void f(const char* caller) \{ cout \textless\textless{} ``When called
from'' \textless\textless{} caller \textless\textless{} ``,''
\textless\textless{} v() \textless\textless{} " gets called.\n``; \} \};
class derived : public base \{ public: derived() \{ f(''derived
constructor``); \} \textasciitilde derived() \{ f(''derived
destructor``); \} const char* v() override \{ return''derived::v()"; \}
\}; int main() \{ derived d; \} Output: When called from base
constructor, base::v() gets called. When called from derived
constructor, derived::v() gets called. When called from derived
destructor, derived::v() gets called. When called from base destructor,
base::v() gets called. The reasoning behind this is that the derived
class may dene additional members which are not yet initialized (in the
constructor case) or already destroyed (in the destructor case), and
calling its member functions would be unsafe. Therefore during
construction and destruction of C++ objects, the dynamic type of
\emph{this is considered to be the constructor's or destructor's class
and not a more-derived class. Example: \#include \#include using
namespace std; class base \{ public: base() \{ std::cout
\textless\textless{} ``foo is'' \textless\textless{} foo()
\textless\textless{} std::endl; \} virtual int foo() \{ return 42; \}
\}; class derived : public base \{ unique\_ptr ptr\_; public:
derived(int i) : ptr\_(new int(i}i)) \{ \} // The following cannot be
called before derived::derived due to how C++ behaves, // if it was
possible\ldots{} Kaboom! int foo() override \{ return \emph{ptr\_; \}
\}; int main() \{ derived d(4); \} Section 38.5: Pure virtual functions
We can also specify that a virtual function is pure virtual (abstract),
by appending = 0 to the declaration. Classes with one or more pure
virtual functions are considered to be abstract, and cannot be
instantiated; only derived classes which dene, or inherit denitions for,
all pure virtual functions can be instantiated. struct Abstract \{
virtual void f() = 0; \}; struct Concrete \{ void f() override \{\} \};
Abstract a; // Error. Concrete c; // Good. Even if a function is specied
as pure virtual, it can be given a default implementation. Despite this,
the function will still be considered abstract, and derived classes will
have to dene it before they can be instantiated. In this case, the
derived class' version of the function is even allowed to call the base
class' version. struct DefaultAbstract \{ virtual void f() = 0; \}; void
DefaultAbstract::f() \{\} struct WhyWouldWeDoThis : DefaultAbstract \{
void f() override \{ DefaultAbstract::f(); \} \}; There are a couple of
reasons why we might want to do this: If we want to create a class that
can't itself be instantiated, but doesn't prevent its derived classes
from being instantiated, we can declare the destructor as pure virtual.
Being the destructor, it must be dened anyways, if we want to be able to
deallocate the instance. And as the destructor is most likely already
virtual to prevent memory leaks during polymorphic use, we won't incur
an unnecessary performance hit from declaring another function virtual.
This can be useful when making interfaces. struct Interface \{ virtual
\textasciitilde Interface() = 0; \};
Interface::\textasciitilde Interface() = default; struct Implementation
: Interface \{\}; // \textasciitilde Implementation() is automatically
defined by the compiler if not explicitly // specified, meeting the
``must be defined before instantiation'' requirement. If most or all
implementations of the pure virtual function will contain duplicate
code, that code can instead be moved to the base class version, making
the code easier to maintain. class SharedBase \{ State my\_state;
std::unique\_ptr my\_helper; // \ldots{} public: virtual void
config(const Context\& cont) = 0; // \ldots{} \}; /} virtual \emph{/
void SharedBase::config(const Context\& cont) \{ my\_helper = new
Helper(my\_state, cont.relevant\_field); do\_this(); and\_that(); \}
class OneImplementation : public SharedBase \{ int i; // \ldots{}
public: void config(const Context\& cont) override; // \ldots{} \}; void
OneImplementation::config(const Context\& cont) /} override */ \{
my\_state = \{ cont.some\_field, cont.another\_field, i \};
SharedBase::config(cont); my\_unique\_setup(); \};  // And so on, for
other classes derived from SharedBase.

\hypertarget{professional-notes-chapter-40-special-member-functions-1}{%
\section{Professional Notes: Chapter 40: Special Member
Functions}\label{professional-notes-chapter-40-special-member-functions-1}}

Section 40.1: Default Constructor A default constructor is a type of
constructor that requires no parameters when called. It is named after
the type it constructs and is a member function of it (as all
constructors are). class C\{ int i; public: // the default constructor
definition C() : i(0)\{ // member initializer list -- initialize i to 0
// constructor function body -- can do more complex things here \} \}; C
c1; // calls default constructor of C to create object c1 C c2 = C(); //
calls default constructor explicitly C c3(); // ERROR: this intuitive
version is not possible due to ``most vexing parse'' C c4\{\}; // but in
C++11 \{\} CAN be used in a similar way C c5{[}2{]}; // calls default
constructor for both array elements C* c6 = new C{[}2{]}; // calls
default constructor for both array elements Another way to satisfy the
``no parameters'' requirement is for the developer to provide default
values for all parameters: class D\{ int i; int j; public: // also a
default constructor (can be called with no parameters) D( int i = 0, int
j = 42 ) : i(i), j(j)\{ \} \}; D d; // calls constructor of D with the
provided default values for the parameters Under some circumstances
(i.e., the developer provides no constructors and there are no other
disqualifying conditions), the compiler implicitly provides an empty
default constructor: class C\{ std::string s; // note: members need to
be default constructible themselves \}; C c1; // will succeed -- C has
an implicitly defined default constructor Having some other type of
constructor is one of the disqualifying conditions mentioned earlier:
class C\{ int i; public: C( int i ) : i(i)\{\} \}; C c1; // Compile
ERROR: C has no (implicitly defined) default constructor Version
\textless{} c++11 To prevent implicit default constructor creation, a
common technique is to declare it as private (with no denition). The
intention is to cause a compile error when someone tries to use the
constructor (this either results in an Access to private error or a
linker error, depending on the compiler). To be sure a default
constructor (functionally similar to the implicit one) is dened, a
developer could write an empty one explicitly. Version c++11 In C++11, a
developer can also use the delete keyword to prevent the compiler from
providing a default constructor. class C\{ int i; public: // default
constructor is explicitly deleted C() = delete; \}; C c1; // Compile
ERROR: C has its default constructor deleted Furthermore, a developer
may also be explicit about wanting the compiler to provide a default
constructor. class C\{ int i; public: // does have automatically
generated default constructor (same as implicit one) C() = default; C(
int i ) : i(i)\{\} \}; C c1; // default constructed C c2( 1 ); //
constructed with the int taking constructor Version c++14 You can
determine whether a type has a default constructor (or is a primitive
type) using std::is\_default\_constructible from : class C1\{ \}; class
C2\{ public: C2()\{\} \}; class C3\{ public: C3(int)\{\} \}; using
std::cout; using std::boolalpha; using std::endl; using
std::is\_default\_constructible; cout \textless\textless{} boolalpha
\textless\textless{} is\_default\_constructible() \textless\textless{}
endl; // prints true cout \textless\textless{} boolalpha
\textless\textless{} is\_default\_constructible() \textless\textless{}
endl; // prints true cout \textless\textless{} boolalpha
\textless\textless{} is\_default\_constructible() \textless\textless{}
endl; // prints true cout \textless\textless{} boolalpha
\textless\textless{} is\_default\_constructible() \textless\textless{}
endl; // prints false Version = c++11 In C++11, it is still possible to
use the non-functor version of std::is\_default\_constructible: cout
\textless\textless{} boolalpha \textless\textless{}
is\_default\_constructible::value \textless\textless{} endl; // prints
true Section 40.2: Destructor A destructor is a function without
arguments that is called when a user-dened object is about to be
destroyed. It is named after the type it destructs with a
\textasciitilde{} prex. class C\{ int* is; string s; public: C() : is(
new int{[}10{]} )\{ \} \textasciitilde C()\{ // destructor definition
delete{[}{]} is; \} \}; class C\_child : public C\{ string s\_ch;
public: C\_child()\{\} \textasciitilde C\_child()\{\} // child
destructor \}; void f()\{ C c1; // calls default constructor C
c2{[}2{]}; // calls default constructor for both elements C* c3 = new
C{[}2{]}; // calls default constructor for both array elements C\_child
c\_ch; // when destructed calls destructor of s\_ch and of C base (and
in turn s) delete{[}{]} c3; // calls destructors on c3{[}0{]} and
c3{[}1{]} \} // automatic variables are destroyed here -- i.e.~c1, c2
and c\_ch Under most circumstances (i.e., a user provides no destructor,
and there are no other disqualifying conditions), the compiler provides
a default destructor implicitly: class C\{ int i; string s; \}; void
f()\{ C* c1 = new C; delete c1; // C has a destructor \} class C\{ int
m; private: \textasciitilde C()\{\} // not public destructor! \}; class
C\_container\{ C c; \}; void f()\{  C\_container* c\_cont = new
C\_container; delete c\_cont; // Compile ERROR: C has no accessible
destructor \} Version \textgreater{} c++11 In C++11, a developer can
override this behavior by preventing the compiler from providing a
default destructor. class C\{ int m; public: \textasciitilde C() =
delete; // does NOT have implicit destructor \}; void f\{ C c1; \} //
Compile ERROR: C has no destructor Furthermore, a developer may also be
explicit about wanting the compiler to provide a default destructor.
class C\{ int m; public: \textasciitilde C() = default; // saying
explicitly it does have implicit/empty destructor \}; void f()\{ C c1;
\} // C has a destructor -- c1 properly destroyed Version \textgreater{}
c++11 You can determine whether a type has a destructor (or is a
primitive type) using std::is\_destructible from : class C1\{ \}; class
C2\{ public: \textasciitilde C2() = delete \}; class C3 : public C2\{
\}; using std::cout; using std::boolalpha; using std::endl; using
std::is\_destructible; cout \textless\textless{} boolalpha
\textless\textless{} is\_destructible() \textless\textless{} endl; //
prints true cout \textless\textless{} boolalpha \textless\textless{}
is\_destructible() \textless\textless{} endl; // prints true cout
\textless\textless{} boolalpha \textless\textless{} is\_destructible()
\textless\textless{} endl; // prints false cout \textless\textless{}
boolalpha \textless\textless{} is\_destructible() \textless\textless{}
endl; // prints false Section 40.3: Copy and swap If you're writing a
class that manages resources, you need to implement all the special
member functions (see Rule of Three/Five/Zero). The most direct approach
to writing the copy constructor and assignment operator would be:
person(const person \&other) : name(new char{[}std::strlen(other.name) +
1{]}) , age(other.age) \{ std::strcpy(name, other.name); \} person\&
operator=(person const\& rhs) \{ if (this != \&other) \{ delete {[}{]}
name;  name = new char{[}std::strlen(other.name) + 1{]};
std::strcpy(name, other.name); age = other.age; \} return \emph{this; \}
But this approach has some problems. It fails the strong exception
guarantee - if new{[}{]} throws, we've already cleared the resources
owned by this and cannot recover. We're duplicating a lot of the logic
of copy construction in copy assignment. And we have to remember the
self-assignment check, which usually just adds overhead to the copy
operation, but is still critical. To satisfy the strong exception
guarantee and avoid code duplication (double so with the subsequent move
assignment operator), we can use the copy-and-swap idiom: class person
\{ char} name; int age; public: /* all the other functions \ldots{}
\emph{/ friend void swap(person\& lhs, person\& rhs) \{ using std::swap;
// enable ADL swap(lhs.name, rhs.name); swap(lhs.age, rhs.age); \}
person\& operator=(person rhs) \{ swap(}this, rhs); return \emph{this;
\} \}; Why does this work? Consider what happens when we have person p1
= \ldots; person p2 = \ldots; p1 = p2; First, we copy-construct rhs from
p2 (which we didn't have to duplicate here). If that operation throws,
we don't do anything in operator= and p1 remains untouched. Next, we
swap the members between }this and rhs, and then rhs goes out of scope.
When operator=, that implicitly cleans the original resources of this
(via the destructor, which we didn't have to duplicate). Self-assignment
works too - it's less ecient with copy-and-swap (involves an extra
allocation and deallocation), but if that's the unlikely scenario, we
don't slow down the typical use case to account for it. Version C++11
The above formulation works as-is already for move assignment. p1 =
std::move(p2); Here, we move-construct rhs from p2, and all the rest is
just as valid. If a class is movable but not copyable, there is no need
to delete the copy-assignment, since this assignment operator will
simply be ill-formed due to the deleted copy constructor. Section 40.4:
Implicit Move and Copy Bear in mind that declaring a destructor inhibits
the compiler from generating implicit move constructors and move
assignment operators. If you declare a destructor, remember to also add
appropriate denitions for the move operations. Furthermore, declaring
move operations will suppress the generation of copy operations, so
these should also be added (if the objects of this class are required to
have copy semantics). class Movable \{ public: virtual
\textasciitilde Movable() noexcept = default; // compiler won't generate
these unless we tell it to // because we declared a destructor
Movable(Movable\&\&) noexcept = default; Movable\&
operator=(Movable\&\&) noexcept = default; // declaring move operations
will suppress generation // of copy operations unless we explicitly
re-enable them Movable(const Movable\&) = default; Movable\&
operator=(const Movable\&) = default; \};
